
\renewcommand{\thetable}{R\arabic{table}}
\renewcommand\thefigure{S\arabic{figure}}

\appendix
\section*{Appendix}

\section{Internal Video Diffusion Model for Research Purpose} \label{subsec:videobackbone}

Our model is a transformer-based latent diffusion model, as illustrated in the~\figref{fig:dit_backbone}. Initially, we employ a 3D VAE to transform videos from the pixel level to a latent space, upon which we construct a transformer-based video diffusion model~\citep{peebles2023scalable}. Previous models, which rely on UNets~\citep{blattmann2023stable,chen2023videocrafter1,guo2023animatediff} or transformers~\citep{ma2024latte}, typically incorporate an additional 1D temporal attention module for video generation, and such spatial-temporally separated designs do not yield optimal results. Instead, we replace the 1D temporal attention with 3D self-attention~\citep{gupta2023photorealistic}, enabling the model to more effectively perceive and process spatiotemporal tokens, thereby achieving a high-quality and coherent video generation model. Specifically, we map the timestep to a scale, thereby applying RMSNorm to the spatiotemporal tokens before each attention or feed-forward network (FFN) module.

\begin{figure*}[ht]
\centering
\includegraphics[width=.9\linewidth]{figs/supp/DiT_backbone.pdf}
\caption{\textbf{Our Video Latent Diffusion Model Backbone}}
\label{fig:dit_backbone}
\end{figure*}


\section{Additional Related Work}

\noindent \textbf{Injecting Control into Video Foundation Models.} \textit{(1) Learning-based}: The control signals are typically projected into latent embeddings via an extra encoder (\eg, learnable convolutional/linear/attention/LoRA layers, or frozen pre-trained feature encoder), which are then integrated into the base model architecture through concatenation, addition, or insertion. VideoComposer[1] employs a unified STC-encoder and CLIP model to feed multi-modal input conditions (textual, spatial, and temporal) into the base T2V model. MotionCtrl~\citep{wang2024motionctrl} introduces camera motion by fine-tuning specific layers of the base U-Net, and object motion via additional convolutional layers. CameraCtrl~\citep{he2024cameractrl} enhances this approach by incorporating ControlNet~\citep{zhang2023adding}'s philosophy, using an attention-based pose encoder to fuse camera signals in the form of Plücker embeddings while keeping the base model frozen. Similarly, SparseCtrl~\citep{guo2023sparsectrl} learns an add-on encoder to integrate control signals (RGB, sketch, depth) into the base model. Tora~\citep{zhang2024tora} employs a trajectory encoder and plug-and-play motion fuser to merge 2D trajectories with the base video model. MotionDirector[7] leverages spatial and temporal LoRA layers to learn desired motion patterns from reference videos.
\textit{(2) Training-free}: These methods modify attention layers or video latents to adjust control signals in a computationally efficient manner. However, training-free methods often suffer from poor generalization and require extensive trial-and-error. Direct-a-video~\citep{yang2024direct} amplifies or suppresses attention in spatial cross-attention layers to inject box guidance, while FreeTraj~\citep{qiu2024freetraj} embeds target trajectories into the low-frequency components and redesigns reweighting strategies across attention layers. MOFT~\citep{xiao2024video} extracts motion priors by removing content correlation and applying motion channel filtering, and then alters the sampling process using the reference MOFT. 

\section{Additional Applications}

We outline our potential applications in various areas as follows.

1) \textbf{Film:} Reproduce the character's classic moves. We can extract the human poses from a given video and apply them to different entities and backgrounds using the capabilities of our model.

2) \textbf{Autonomous Driving:}  Simulate dangerous safety accidents, such as two cars colliding and a car hitting a person.

3) \textbf{Embodied AI:} Generate a vast number of videos with diverse entity and trajectory inputs to train a general 4D pose estimator, especially for non-rigid objects.

4) \textbf{Game:} Train a character ID, such as Black Myth Wukong, through LoRA, and then drive the character movement with different trajectories.

\section{Clarification of the Limited Entity Number ($\leq$3)} 

Currently, our method is limited to generating up to 3 entities, as outlined in the `Limitation' section of the paper. This constraint is primarily due to the capabilities of the video foundation model rather than the training data. While it is relatively easy to generate $\gg$2 entities of the same category (e.g., ``a group of people/cars/animals") in the video, it becomes much more challenging to generate $\gg$2 entities, each differing greatly from the others, through the text input as T5 text encoder tends to mix the textual features of different entities. Thus it becomes hard to associate specific trajectories with their corresponding text entities. Based on empirical studies with video foundation models, we chose to limit the number of entities to 3 in our work. Regarding data construction, it is easy to include more entities with their paired trajectories in our procedure UE platform pipeline. However, the key limitation is that the video foundation model struggles to generate such a diverse set of entities simultaneously. Furthermore, many prior works, such as Tora, MotionCtrl, and Direct-a-video also focus on a limited number of entities.

\section{Dataset Illustration} 

\subsection{360$^{\circ}$-Motion Dataset Data.} \label{subsec:360motiondataset}
We show a sample in~\figref{fig:supp_dataset} captured with 12 evenly-surrounded cameras. Each camera shoots a clip of 100 frames at 384$\times$672 resolutions. During training, we discard the initial 10 frames to eliminate potential blurring and noise caused by 3D model initialization in the UE platform.

\begin{figure*}[!ht]
\centering
\includegraphics[width=\linewidth]{figs/supp/supp_dataset.pdf}
\caption{\textbf{A Sample from our \textit{360$^{\circ}$-Motion} Captured with 12 Evenly-Surrounded Cameras.}}
\label{fig:supp_dataset}
\end{figure*}

\subsection{Unbalanced Entity Distribution in Common Video Datasets} \label{subsec:unbalanceddataset}

In high-quality video datasets like Artgrid, Pixabay, and Pexels\footnote{Artgrid: https://artgrid.io/, Pixabay: https://www.videvo.net/, Pexels: https://www.pexels.com/}, the issue of category imbalance is highly pronounced and poses significant challenges. We analyze the aforementioned three datasets by first captioning the videos using QWen-VL~\citep{Qwen-VL}. Subsequently, we employ the spaCy\footnote{spaCy: https://spacy.io/} library to extract noun chunks from the video captions, which serve as entity words. We predefine over 60 classes as keywords for entity filtering. As illustrated in the~\figref{fig:entity_distribution}, certain categories (\eg, humans) constitute a disproportionately large share of the entity objects, thereby constraining the model's ability to generalize to other categories that appear less frequently.

\begin{figure*}[!ht]
\centering
\includegraphics[width=\linewidth]{figs/supp/dataset_distribution.png}
\caption{\textbf{Entity Distribution Over 60 Classes in Artgrid, Pixabay, and Pexels.}}
\label{fig:entity_distribution}
\end{figure*}

\subsection{GPT-Generated Evaluation Prompts} \label{subsec:evaluateprompt}

The human prompts, non-human (animal, car, robot) prompts, and location prompts for evaluation are provided in~\tabref{tab:human_prompts}, ~\tabref{tab:nonhuman_prompts_1}\&\tabref{tab:nonhuman_prompts_2}, and~\tabref{tab:location_prompts} respectively.

\begin{table}[!ht]

\caption{\textbf{Evaluation Human Prompts}. They are generated using GPT prompt: \textit{``Generate more human samples similar to \{Train Human Sample\}, no more than 25 words.''}}
\label{tab:human_prompts}

\begin{threeparttable}
\resizebox{.99\textwidth}{!} 
{
\begin{tabular}{l}
\toprule

1. a man with short spiky brown hair, athletic build, a navy blue jacket, beige cargo pants, \\ \quad and black sneakers  \\
2. a woman with long wavy blonde hair, petite figure, a red floral dress, white sandals, and a \\ \quad yellow shoulder bag \\
3. a man with a shaved head, broad shoulders, a gray graphic t-shirt, dark jeans, and brown \\ \quad leather boots \\
4. a woman with shoulder-length straight auburn hair, a slender figure, a green button-up \\ \quad 
 blouse, black leggings, and white sneakers \\
5. a man with messy black hair, tall frame, a plaid red and black shirt, faded blue jeans, and \\ \quad tan hiking boots \\
6. a man with medium-length straight brown hair, tall and slender, a gray crew-neck t-shirt,\\ \quad beige trousers, and dark green sneakers \\
7. a woman with short curly black hair, slender build, a pink hoodie, light gray joggers, and \\ \quad blue sneakers \\
8. a man with short black wavy hair, lean figure, a green and yellow plaid shirt, dark brown \\ \quad pants, and black suede shoes \\
9. a man with curly black hair, muscular build, a dark green hoodie, gray joggers, and white \\ \quad running shoes \\
10. a woman with short blonde hair, slim athletic build, a red leather jacket, dark blue jeans, \\ \quad and white sneakers \\
11. a man with medium-length wavy brown hair, lean build, a black bomber jacket, olive \\ \quad green cargo pants, and brown hiking boots \\
12. a man with buzz-cut blonde hair, stocky build, a gray zip-up sweater, black shorts, and red \\ \quad basketball shoes \\
13. a woman with long straight black hair, toned build, a blue denim jacket, light gray legg\\ \quad -ings, and black slip-on shoes \\
14. a man with short curly red hair, average build, a black leather jacket, dark blue cargo pants, \\ \quad and white sneakers \\
15. a woman with shoulder-length wavy brown hair, slim build, a green parka, black leggings, \\ \quad and gray hiking boots \\
16. a man with short straight black hair, tall and lean build, a navy blue sweater, khaki shorts, \\ \quad and brown sandals \\
17. a woman with pixie-cut blonde hair, athletic build, a red windbreaker, blue ripped jeans, \\ \quad and black combat boots \\
18. a man with medium-length wavy gray hair, muscular build, a maroon t-shirt, beige chinos, \\ \quad and brown loafers \\
19. a woman with long curly black hair, average build, a purple hoodie, black athletic shorts, \\ \quad and white running shoes \\
20. a man with short spiky blonde hair, slim build, a black trench coat, blue jeans, and brown \\ \quad hiking shoes \\

\bottomrule

\end{tabular}
}

\end{threeparttable}
\end{table}

\begin{table}[!ht]

\caption{\textbf{Evaluation Non-Human Prompts (1/2)}. They are generated using GPT prompt: \textit{``Generate more animal/car/robot samples similar to \{Train Sample\}, no more than 25 words.''}}
\label{tab:nonhuman_prompts_1}

\begin{threeparttable}
\resizebox{.99\textwidth}{!} 
{
\begin{tabular}{l}
\toprule
1. a dog with a fluffy coat, wagging tail, and warm golden-brown fur, exuding a gentle and \\ \quad friendly charm \\
2. a tiger with vibrant orange and black stripes, piercing yellow eyes, and a powerful stance, \\ \quad exuding strength and grace \\
3. a giraffe with golden-yellow fur, long legs, a tall slender neck, and patches of brown spots, \\ \quad exuding elegance and calm \\
4. an alpaca with soft white wool, short legs, a thick neck, and a fluffy head of fur, radiating \\ \quad gentle charm \\
5. a zebra with black and white stripes, sturdy legs, a short neck, and a sleek mane running \\ \quad down its back \\
6. a deer with sleek tan fur, long slender legs, a graceful neck, and tiny antlers atop its head \\
7. a gazelle with light golden fur, long slender legs, a thin neck, and short, sharp horns, \\ \quad embodying elegance and agility \\
8. a horse with chestnut brown fur, muscular legs, a slim neck, and a flowing mane, exuding \\ \quad strength and grace \\
9. a sleek black panther with a smooth, glossy coat, emerald green eyes, and a powerful stance \\
10. a cheetah with golden fur covered in black spots, intense amber eyes, and a slender, \\ \quad agile body \\
11. a regal lion with a thick, flowing golden mane, sharp brown eyes, and a powerful muscular \\ \quad frame \\
12. a snow leopard with pale gray fur adorned with dark rosettes, icy blue eyes, and a stealthy, \\ \quad poised posture \\
13. a jaguar with a golden-yellow coat dotted with intricate black rosettes, deep green eyes, \\ \quad and a muscular build \\
14. a wolf with thick silver-gray fur, alert golden eyes, and a lean yet strong body, exuding \\ \quad confidence and boldness \\
15. a tiger with a pristine white coat marked by bold black stripes, bright blue eyes, and a \\ \quad graceful, poised form \\
16. a lynx with tufted ears, soft reddish-brown fur with faint spots, and intense yellow-green \\ \quad eyes \\
17. a bear with dark brown fur, small but fierce black eyes, and a broad and muscular build, \\ \quad radiating power \\
18. a swift fox with reddish-orange fur, a bushy tail tipped with white, and sharp, intelligent \\ \quad amber eyes \\
19. a falcon with blue-gray feathers, sharp talons, and keen yellow eyes fixed on its prey below \\
20. a fox with sleek russet fur, a bushy tail tipped with black, and bright green and cunning eyes \\
21. a kangaroo with brown fur, powerful hind legs, and a muscular tail, showcasing its strength \\ \quad and agility \\
22. a polar bear with thick white fur, strong paws, and a black nose, embodying the essence \\ \quad of the Arctic \\
23. a cheetah with a slender build, spotted golden fur, and sharp eyes, epitomizing speed and \\ \quad agility \\
24. a dolphin with sleek grey skin, a curved dorsal fin, and intelligent, playful eyes, reflecting \\ \quad its nature \\
25. a wolf with a body covered in thick silver fur, sharp ears, and piercing yellow eyes, \\ \quad showcasing its alertness \\
26. a leopard with a body covered in golden fur, dark rosettes, and a long muscular tail, \\ \quad emphasizing its strength \\
27. a penguin with a body covered in smooth black-and-white feathers, short wings, and \\ \quad webbed feet \\
28. a gazelle with a body covered in sleek tan fur, long legs, and elegant curved horns, \\showcasing its grace \\

\bottomrule

\end{tabular}
}

\end{threeparttable}
\end{table}

\begin{table}[!ht]

\caption{\textbf{Evaluation Non-Human Prompts (2/2)}. They are generated using GPT prompt: \textit{``Generate more animal/car/robot samples similar to \{Train Sample\}, no more than 25 words.''}}
\label{tab:nonhuman_prompts_2}

\begin{threeparttable}
\resizebox{.99\textwidth}{!} 
{
\begin{tabular}{l}
\toprule
29. a rabbit with a body covered in soft fur, quick hops, and a playful demeanor, showcasing \\ \quad its energy \\
30. a koala with a body covered in soft grey fur, large round ears, and a black nose, radiating \\ \quad cuteness \\
31. a rhinoceros with a body covered in thick grey skin, a massive horn on its snout, and \\ \quad sturdy legs \\
32. a flamingo with a body covered in pink feathers, long slender legs, and a gracefully \\ \quad curved neck \\
33. a parrot with bright red, blue, and yellow feathers, a curved beak, and sharp eyes \\
34. a hippopotamus with a body covered in thick grey-brown skin, massive jaws, and a \\ \quad large body \\
35. a crocodile with a body covered in scaly green skin, a powerful tail, and sharp teeth \\
36. a moose with a body covered in thick brown fur, massive antlers, and a bulky frame \\
37. a fluttering butterfly with intricate wing patterns, vivid colors, and graceful flight \\
38. a chameleon with a body covered in vibrant green scales, bulging eyes, and a curled tail, \\ \quad showcasing its unique charm \\
39. a lemur with a body covered in soft grey fur, a ringed tail, and wide yellow eyes, and \\ \quad curious expression \\
40. a squirrel with a body covered in bushy red fur, large eyes, and a fluffy tail \\
41. a panda with a body covered in fluffy black-and-white fur, a round face, and gentle eyes, \\ \quad radiating warmth \\
42. a porcupine with a body covered in spiky brown quills, a small nose, and curious eyes \\
43. a sedan with a sleek metallic silver body, long wheelbase, a low-profile hood, and a small \\ \quad rear spoiler \\
44. an SUV with a matte black exterior, elevated suspension, a tall roofline, and a compact \\ \quad rear roof rack \\
45. a pickup truck with rugged dark green paint, extended cab, raised suspension, and a modest \\ \quad cargo bed cover \\
46. a vintage convertible with a body covered in shiny red paint, chrome bumpers, and a stylish \\ \quad design \\
47. a futuristic electric car with a minimalist silver design, slim LED lights, and smooth curves \\
48. a compact electric vehicle with a silver finish, aerodynamic profile, and efficient battery \\
49. a firefighting robot with a water cannon arm, heat sensors, and durable red-and-silver exterior \\
50. an industrial welding robot with articulated arms, a laser precision welder, and heat-resistant \\ \quad shields \\
51. a disaster rescue robot with reinforced limbs, advanced AI, and a rugged body designed \\ \quad to navigate \\
52. an exploration rover robot with solar panels, durable wheels, and advanced sensors for \\ \quad planetary exploration \\

\bottomrule

\end{tabular}
}

\end{threeparttable}
\end{table}

\begin{table}[!ht]

\caption{\textbf{Evaluation Location Prompts}. }
\label{tab:location_prompts}

\begin{threeparttable}
\resizebox{.99\textwidth}{!} 
{
\begin{tabular}{l}
\toprule
1. fjord 2. sunset beach 3. cave 4. snowy tundra 5. prairie 6. asian town 7. rainforest 8. canyon\\
9. savanna 10. urban rooftop garden 11. swamp 12. riverbank 13. coral reef 14. volcanic landscape \\
15. wind farm 16. town street 17. night city square 18. mall lobby 19. glacier 20. seaside street 
\\
21. gymnastics room 22. abandoned factory 23. autumn forest 24. mountain village 25. coastal harbor \\
26. ancient ruins 27. modern metropolis 28. dessert 29. forest 30. city 31. snowy street 32. park \\

\bottomrule

\end{tabular}
}

\end{threeparttable}
\end{table}


\section{More Experiments}

\subsection{Fine-grained Entity Prompt Input}

We provide additional samples in~\figref{fig:supp_diffhuman} to demonstrate that 3DTrajMaster supports fine-grained entity customization. The description of the man can be flexibly modified by adjusting attributes such as hair, gender, physique, clothing, and accessories.

\begin{figure*}[!ht]
\centering
\includegraphics[width=\linewidth]{figs/supp/supp_diffhuman.pdf}
\caption{\textbf{Flexible Entity Editing in Input Text Prompts.} The other entity, \textit{``a swift falcon with blue-gray feathers, sharp talons, and keen yellow eyes focused on its prey below"} remains fixed while varying the human entity descriptions.}
\label{fig:supp_diffhuman}
\end{figure*}

\subsection{Ablation Study}

\subsubsection{Optimal Hyperparameters} \label{subsec:optimialhyper}

In the main paper, we propose a video domain adaptor and an annealed sampling strategy to mitigate video domain shifts from our constructed UE datasets. However, completely removing the LoRA adaptor (as the learned motion and domain bias are coupled to some extent) or the inserted motion guidance will result in a decline in 3D trajectory accuracy. Thus, applying video enhancement techniques with appropriate dropping is crucial. To this end, we begin with a randomly initialized parameter group: $T_{c}=10, \alpha=0.2, TS=72,000$. We perform ablation experiments on our evaluation subset. As shown in~\tabref{tab:ablation_annealed_tc},~\tabref{tab:ablation_lora_alpha}, and~\tabref{tab:ablation_ts}, the video quality exhibits a monotonically decreasing trend as these hyperparameters increase. In contrast, 3D trajectory accuracy initially drops sharply but stabilizes in the later stages. To balance the degradation of visual quality with maintaining pose accuracy, we select an optimal parameter group: $T_{c}=25, \alpha=0.4, TS=36,000$ as our default inference setting.


\begin{table}[!ht]

\caption{Ablation Study on Annealed Timestep $T_c$.}
\label{tab:ablation_annealed_tc}

\centering
\begin{threeparttable}

\resizebox{0.85\textwidth}{!} 
{
\begin{tabular}{ccccccccccc}
\toprule
&  \multicolumn{3}{c}{ Video Quality } & \multicolumn{2}{c}{ 3D Trajectory Accuracy } \\
\cmidrule(lr){2-4} 
\cmidrule(lr){5-6} 
Annealed Timestep $T_c$ & FVD  & FID  & CLIPSIM $\uparrow$ & TransErr (m)  & RotErr (deg)  \\

\midrule
$T_{c}=5$ & \textbf{1492.79} & \textbf{76.95} & \textbf{0.3469} & 0.844 & 1.099 \\
$T_{c}=10$ & \underline{1976.01} & \underline{106.45} & \underline{0.3429} & 0.546 & 0.493  \\
$T_{c}=15$ & 2179.15 & 122.55 & 0.3405 & 0.437 & 0.422 \\
$T_{c}=20$ & 2236.05 & 128.89 & 0.3374 & 0.391 & 0.284 \\
$T_{c}=25$ & 2240.40 & 132.90 & 0.3337 & \textbf{0.344} & 0.274 \\
$T_{c}=30$ & 2295.13 & 137.52 & 0.3314 & 0.360 & \textbf{0.261} \\
$T_{c}=35$ & 2323.20 & 142.71 & 0.3276 & 0.352 & \underline{0.264} \\
$T_{c}=40$ & 2338.47 & 148.27 & 0.3240 & 0.351 & 0.266 \\
$T_{c}=45$ & 2363.49 & 156.39 & 0.3207 & 0.350 & 0.268 \\
$T_{c}=50$ & 2347.64 & 166.71 & 0.3185 & \underline{0.348} & 0.281 \\

\bottomrule

\end{tabular}
} 
\end{threeparttable}
\end{table}

\begin{table}[!ht]

\caption{Ablation Study on LoRA Scalar $\alpha$.} 
\label{tab:ablation_lora_alpha}

\centering
\begin{threeparttable}
\resizebox{.85\textwidth}{!} 
{
\begin{tabular}{ccccccccccc}
\toprule
&  \multicolumn{3}{c}{ Video Quality } & \multicolumn{2}{c}{ 3D Trajectory Accuracy } \\
\cmidrule(lr){2-4} 
\cmidrule(lr){5-6} 
LoRA Scalar $\alpha$ & FVD  & FID  & CLIPSIM $\uparrow$ & TransErr (m)  & RotErr (deg)    \\

\midrule
$\alpha=0$ & \textbf{1495.38} & \textbf{80.56} & \textbf{0.3467} & 0.646 & 0.900  \\
$\alpha=0.2$ & \underline{1976.01} & \underline{106.45} & \underline{0.3429}  & 0.546 & 0.493 \\
$\alpha=0.4$ & 2150.42 & 133.76 & 0.3367 & 0.444 & \underline{0.428}   \\
$\alpha=0.6$ & 2330.56 & 152.12 & 0.3277 & 0.394 & \textbf{0.393}  \\
$\alpha=0.8$ & 2318.78 & 195.93 & 0.3125 & \underline{0.378} & 0.450  \\
$\alpha=1.0$ & 2481.33 & 224.81 & 0.3087 & \textbf{0.358} & 0.432 \\
\bottomrule


\end{tabular}
}

\end{threeparttable}
\end{table}

\begin{table}[!ht]

\caption{Ablation Study on Training Step $TS$.}
\label{tab:ablation_ts}

\centering
\begin{threeparttable}
\resizebox{.85\textwidth}{!} 
{
\begin{tabular}{ccccccccccc}
\toprule
&  \multicolumn{3}{c}{ Video Quality } & \multicolumn{2}{c}{ 3D Trajectory Accuracy } \\
\cmidrule(lr){2-4} 
\cmidrule(lr){5-6} 
Train. Steps $TS$ & FVD  & FID  & CLIPSIM $\uparrow$ & TransErr (m)  & RotErr (deg)  \\

\midrule
$TS=12,000$ & \textbf{1493.68} & \textbf{72.03} & \underline{0.3427} & 0.561 & 0.713 \\
$TS=36,000$ & \underline{1883.15} & \underline{99.98} & 0.3408 & 0.523 & 0.631 \\
$TS=72,000$ & 1976.01 & 106.45 & \textbf{0.3429} & 0.546 & 0.493 \\
$TS=108,000$ & 2068.43 & 111.01 & 0.3388 & \underline{0.446} & \textbf{0.480} \\
$TS=144,000$ & 2102.28 & 114.84 & 0.3367 & \textbf{0.411} & \underline{0.482} \\
\bottomrule

\end{tabular}
}

\end{threeparttable}
\end{table}


\subsubsection{Negative Pose Condition as Static Motions} \label{subsec:negpose}

We find that setting negative pose sequences as static motions $\{(\hat{\mathbf{P}}_n)_{n=1}^N|\hat{\mathbf{P}}_n=\mathbf{P}_0,\forall n\}$ rather than positive motion sequences $\{(\mathbf{P}_n)_{n=1}^N\}$ can further improve pose accuracy, as shown in~\tabref{tab:ablation_neg_pose}. We infer that the model captures underlying 3D motion representations from the randomly generated 3D trajectories. However, we do not adopt this approach due to the decline in video quality.

\begin{table}[!ht]

\caption{Ablation Study on Negative Pose Sequences.}
\label{tab:ablation_neg_pose}

\centering
\begin{threeparttable}
\resizebox{.92\textwidth}{!} 
{
\begin{tabular}{ccccccccccc}
\toprule
&  \multicolumn{3}{c}{ Video Quality } & \multicolumn{2}{c}{ 3D Trajectory Accuracy } \\
\cmidrule(lr){2-4} 
\cmidrule(lr){5-6} 
Negative Condition & FVD  & FID  & CLIPSIM $\uparrow$ & TransErr (m)  & RotErr (deg)  \\

\midrule
Neg. Pose = Static Motions  & 2141.39 & 118.22 & 0.3360 & \textbf{0.371} & \textbf{0.448} \\
Neg. Pose = Pos. Pose  & \textbf{1976.01} & \textbf{106.45} & \textbf{0.3429} & 0.546 & 0.493 \\
\bottomrule

\end{tabular}
}

\end{threeparttable}
\end{table}


\subsubsection{Qualitative Feedback from Human Users}

We conducted a questionnaire survey and collected 53 samples to form user preference comparisons. Each participant received a reward of 0.80 USD and spent approximately 5 minutes completing the questionnaire, which assessed four dimensions: (1) video quality, (2) trajectory accuracy, (3) entity diversity, and (4) background diversity. In~\tabref{tab:user_study}, we report the proportion of users who preferred our model over the baselines.

\begin{table}[!ht]

\caption{User Preference Comparisons.}
\label{tab:user_study}

\centering
\begin{threeparttable}
\resizebox{.57\textwidth}{!} 
{
\begin{tabular}{c|c|c|c}
\toprule
Method & MotionCtrl & Direct-a-Video &  Tora \\
\midrule 3DTrajMaster & $47.2 \%$ & $56.6 \%$ & $81.1 \%$ \\
\bottomrule
\end{tabular}
}

\end{threeparttable}
\end{table}



\subsubsection{Generalizable Entity Prompts\&3D Trajectories} \label{subsec:generalizableresults}

We provide more generalizable results with novel entity prompts generated by GPT and 3D trajectories, as shown in~\figref{fig:sub_generalizable_start} to~\figref{fig:sub_generalizable_end}. Each text prompt consists of one to three entities. \textcolor{red}{\textit{(We kindly urge readers to check the visual results in the our website)}}.

\begin{figure*}[!ht]
\centering
\includegraphics[width=.94\linewidth]{figs/supp/single_entity_1.pdf}
\caption{\textbf{Generalizable Results with Novel 3D Trajectories \& Entity Prompts (1/20)}}
\label{fig:sub_generalizable_start}
\end{figure*}

\begin{figure*}[!ht]
\centering
\includegraphics[width=.93\linewidth]{figs/supp/single_entity_2.pdf}
\caption{\textbf{Generalizable Results with Novel 3D Trajectories \& Entity Prompts (2/20)}}
\end{figure*}

\begin{figure*}[!ht]
\centering
\includegraphics[width=.93\linewidth]{figs/supp/single_entity_3.pdf}
\caption{\textbf{Generalizable Results with Novel 3D Trajectories \& Entity Prompts (3/20)}}
\end{figure*}

\begin{figure*}[!ht]
\centering
\includegraphics[width=.93\linewidth]{figs/supp/single_entity_4.pdf}
\caption{\textbf{Generalizable Results with Novel 3D Trajectories \& Entity Prompts (4/20)}}
\end{figure*}

\begin{figure*}[!ht]
\centering
\includegraphics[width=.93\linewidth]{figs/supp/two_entity_1.pdf}
\caption{\textbf{Generalizable Results with Novel 3D Trajectories \& Entity Prompts (5/20)}}
\end{figure*}

\begin{figure*}[!ht]
\centering
\includegraphics[width=.93\linewidth]{figs/supp/two_entity_2.pdf}
\caption{\textbf{Generalizable Results with Novel 3D Trajectories \& Entity Prompts (6/20)}}
\end{figure*}

\begin{figure*}[!ht]
\centering
\includegraphics[width=.93\linewidth]{figs/supp/two_entity_3.pdf}
\caption{\textbf{Generalizable Results with Novel 3D Trajectories \& Entity Prompts (7/20)}}
\end{figure*}

\begin{figure*}[!ht]
\centering
\includegraphics[width=.93\linewidth]{figs/supp/two_entity_4.pdf}
\caption{\textbf{Generalizable Results with Novel 3D Trajectories \& Entity Prompts (8/20)}}
\end{figure*}

\begin{figure*}[!ht]
\centering
\includegraphics[width=.93\linewidth]{figs/supp/two_entity_5.pdf}
\caption{\textbf{Generalizable Results with Novel 3D Trajectories \& Entity Prompts (9/20)}}
\end{figure*}

\begin{figure*}[!ht]
\centering
\includegraphics[width=.93\linewidth]{figs/supp/two_entity_6.pdf}
\caption{\textbf{Generalizable Results with Novel 3D Trajectories \& Entity Prompts (10/20)}}
\end{figure*}

\begin{figure*}[!ht]
\centering
\includegraphics[width=.93\linewidth]{figs/supp/two_entity_7.pdf}
\caption{\textbf{Generalizable Results with Novel 3D Trajectories \& Entity Prompts (11/20)}}
\end{figure*}

\begin{figure*}[!ht]
\centering
\includegraphics[width=.93\linewidth]{figs/supp/two_entity_8.pdf}
\caption{\textbf{Generalizable Results with Novel 3D Trajectories \& Entity Prompts (12/20)}}
\end{figure*}

\begin{figure*}[!ht]
\centering
\includegraphics[width=.93\linewidth]{figs/supp/two_entity_9.pdf}
\caption{\textbf{Generalizable Results with Novel 3D Trajectories \& Entity Prompts (13/20)}}
\end{figure*}

\begin{figure*}[!ht]
\centering
\includegraphics[width=.93\linewidth]{figs/supp/two_entity_10.pdf}
\caption{\textbf{Generalizable Results with Novel 3D Trajectories \& Entity Prompts (14/20)}}
\end{figure*}

\begin{figure*}[!ht]
\centering
\includegraphics[width=.93\linewidth]{figs/supp/two_entity_11.pdf}
\caption{\textbf{Generalizable Results with Novel 3D Trajectories \& Entity Prompts (15/20)}}
\end{figure*}

\begin{figure*}[!ht]
\centering
\includegraphics[width=.93\linewidth]{figs/supp/three_entity_1.pdf}
\caption{\textbf{Generalizable Results with Novel 3D Trajectories \& Entity Prompts (16/20)}}
\end{figure*}

\begin{figure*}[!ht]
\centering
\includegraphics[width=.93\linewidth]{figs/supp/three_entity_2.pdf}
\caption{\textbf{Generalizable Results with Novel 3D Trajectories \& Entity Prompts (17/20)}}
\end{figure*}

\begin{figure*}[!ht]
\centering
\includegraphics[width=.93\linewidth]{figs/supp/three_entity_3.pdf}
\caption{\textbf{Generalizable Results with Novel 3D Trajectories \& Entity Prompts (18/20)}}
\end{figure*}

\begin{figure*}[!ht]
\centering
\includegraphics[width=.93\linewidth]{figs/supp/three_entity_4.pdf}
\caption{\textbf{Generalizable Results with Novel 3D Trajectories \& Entity Prompts (19/20)}}
\end{figure*}

\begin{figure*}[!ht]
\centering
\includegraphics[width=.93\linewidth]{figs/supp/three_entity_5.pdf}
\caption{\textbf{Generalizable Results with Novel 3D Trajectories \& Entity Prompts (20/20)}}
\label{fig:sub_generalizable_end}
\end{figure*}
